\section{Method}
\label{sec:method}

\subsection{The idea, before the math}

After the first decoder layer fires, the model already has a rough opinion about where objects are. We treat that opinion as new information about the image and use it to refine the encoder memory. Layers $2$ through $L$ then attend over this refined memory, so the decoder reads features that already incorporate its own early hypotheses.

Concretely, after decoder layer $1$ produces its query outputs $\mathbf{t}^{(1)}$, we run a single cross-attention pass with the encoder memory $\mathbf{m}$ as queries and $\mathbf{t}^{(1)}$ as keys and values, and gate the result back into $\mathbf{m}$ with a learnable scalar $g$.

\subsection{Formula 1: attention}

Standard scaled dot-product attention~\cite{vaswani2017attention} reads
\begin{equation}
    \mathrm{Attn}(\mathbf{Q}, \mathbf{K}, \mathbf{V}) \;=\; \mathrm{softmax}\!\left(\tfrac{\mathbf{Q}\mathbf{K}^{\top}}{\sqrt{d}}\right) \mathbf{V}.
    \label{eq:attn}
\end{equation}

\noindent\textbf{Variables.} $\mathbf{Q}$ is what we look for; $\mathbf{K}$ is where we look (the addresses); $\mathbf{V}$ is what we extract (the contents). $d$ is the per-head feature dimension; the scaling keeps the softmax's input variance well-behaved.

\noindent\textbf{Intuition.} Each query computes a similarity score against every key, normalizes those scores into weights, and reads out a weighted average of the values. A query "asks a question of the sequence", and the answer is a soft retrieval over its contents.

\subsection{Formula 2: feedback update}

Our feedback step is one application of attention with a residual gate:
\begin{equation}
    \mathbf{m}' \;=\; \mathbf{m} \;+\; g \cdot \mathrm{Attn}(\mathbf{m}, \mathbf{t}^{(1)}, \mathbf{t}^{(1)}).
    \label{eq:feedback}
\end{equation}

\noindent\textbf{Variables.} $\mathbf{m}$ is the encoder memory --- the image features the decoder reads from; $\mathbf{t}^{(1)}$ is the decoder's output after its first layer (its preliminary predictions); $g \in (0, 1)$ is a learnable scalar gate that controls how much of the refinement is mixed back into memory.

\noindent\textbf{Intuition.} Each memory token asks the early decoder predictions, "is anything you found relevant to me?" The attention answer is added back to the original memory, weighted by $g$. When $g \to 0$ the module is a no-op; when $g \to 1$ the refinement dominates. The optimizer chooses where on this continuum to operate, and the value it picks turns out to be the entire story of this work.

In practice we use multi-head attention, follow Equation~\ref{eq:feedback} with a small feed-forward block, and apply LayerNorm to both residual paths --- standard transformer plumbing that does not change the picture above.
